% Part: proof-theory
% Chapter: normalization
% Section: translations

\documentclass[../../../include/open-logic-section]{subfiles}

\begin{document}

\olfileid{pt}{nor}{tr}

\olsection{સામાન્ય \Log{N2i}-\usetoken{P}{proof} અને \Log{G2i} વચ્ચેનું રૂપાંતર}

!!^{proof}sને
\Log{N2i}માંથી
\Log{G2i}ની
!!{proof}sમાં
રૂપાંતરિત કરી
શકાય છે અને
ઊલટું પણ.
પહેલાં, કટ-રહિત
\Log{G2i}ની
!!{proof}sનું
રૂપાંતર સામાન્ય
\Log{N2i}-!!{proof}s
આપે છે.
\sourcecorrection{OLMD-261}{મૂળમાં કટ-રહિત સાબિતીના તંત્રનું નામ \Log{2i} છે, જે અપરિભાષિત છે; અહીં સંદર્ભ અને પૂર્વવર્તી રૂપાંતર મુજબ \Log{G2i} છે.}

પરિણામ કહેવા
\Log{N2i} અને
\Log{G2i}ના
સિક્વન્ટોના
પૂર્વાંગો વચ્ચેનો
સંબંધ સરળતાથી
દર્શાવતું થોડું
પરિભાષિક ગદ્ય
જોઈએ: ચિહ્નિત
!!{formula}sનો
ગણ~$\Gamma'$
ચિહ્નવિહોણાં
!!{formula}sના
બહુગણ~$\Gamma$ને
\emph{અનુરૂપ છે},
જો દરેક
!!{formula}~$!A$
માટે~$\Gamma'$માં
$x : !A$
સ્વરૂપના
!!{element}sની
સંખ્યા~$\Gamma$માં
$!A$ની બહુલતા
કરતાં ઓછી કે
સમાન હોય.

\begin{cor}
જો
$\Log{G2i} \Proves \Gamma \Sequent \Delta$
હોય, તો
$\Gamma' \Sequent !C$ની
સામાન્ય
\Log{N2i}-!!{proof}~$\delta$
હોય છે. પહેલાના
સિક્વન્ટના જમણા
ભાગમાં સૂત્ર
હોય તો તે
જ પરિણામનું
સૂત્ર છે; જમણો
ભાગ ખાલી હોય
તો પરિણામ
અસત્ય છે.
અહીં
ચિહ્નિત
!!{formula}sનો
ગણ~$\Gamma'$
બહુગણ~$\Gamma$ને
અનુરૂપ છે:
એટલે દરેક
!!{formula}~$!A$
માટે~$\Gamma'$માં
$x : !A$
સ્વરૂપના
!!{element}sની
સંખ્યા~$\Gamma$માં
$!A$ની બહુલતા
(એટલે~$!A$ની
નકલોની સંખ્યા)
કરતાં ઓછી કે
સમાન છે.
\sourcecorrection{OLMD-262}{મૂળ ઉપપત્તિ કટ-રહિત G2iમાંથી N2cની સામાન્ય સાબિતી કહે છે; અગાઉના રૂપાંતર મુજબ લક્ષ્ય \Log{N2i} છે. મૂળમાં જમણો ભાગ ખાલી હોય ત્યારે પરિણામ અસત્ય લેવાય અને એક સૂત્ર હોય ત્યારે તે જ લેવાય તે શરત પણ છૂટી છે.}
\end{cor}

\begin{proof}
\olref[nat][tgi]{prop:G2i-to-N2i}ની
સાબિતી તપાસો:
!!{introduction} નિયમો
!!{proof}sના
અંતે ઉમેરાય છે,
અને !!{elimination}
નિયમો ફક્ત
!!{proof}sની
ટોચે ઉમેરાય
છે. તેથી
!!a{elimination}
નિયમની ઉપર
!!{introduction}
નિયમ કદી
આવતો નથી.
\end{proof}

જોકે
\olref[nat][tgi]{prop:G2i-to-N2i}માં
આપેલું~\Cut{}
નિયમનું રૂપાંતર
\Log{N2i}-!!{proof}~$\delta$માં
કટ ઊભા કરી
શકે છે.
\Log{G2i}ની
!!{proof}નો
અંત~\Cut{}માં
થતો હોય અને
તેની તત્કાલીન
ઉપ-!!{proof}s
$\pi_1$ અને
$\pi_2$ અનુક્રમે
$\Gamma_1
\Sequent !A$
અને
$!A, \Gamma_2 \Sequent \Delta_2$
ની હોય, તો
$\pi_1$નું
રૂપાંતર~$\delta_1$,
$\pi_2$ના
રૂપાંતર~$\delta_2$માં
આવતા
$x:!A
\Sequent !A$
સ્વયંસિદ્ધો પર
કલમ-જોડીએ છીએ.
$\delta_1$નો
અંત મુખ્ય
!!{formula}~$!A$
ધરાવતા
!!a{introduction}
નિયમમાં થાય
અને~$\delta_2$માંનો
$x:!A \Sequent !A$
સ્વયંસિદ્ધ
!!a{elimination}
નિયમનું મુખ્ય
આધારવિધાન હોય,
તો~$\delta$માં
કટ ઊભો થાય.

\Log{N2i}ની
!!{proof}sને
\Log{G2i}ની
!!{proof}sમાં
પણ રૂપાંતરિત
કરી શકાય છે.
\olref[nat][tni]{prop:N2-to-G2}ની
સાબિતીમાં
\Elim{\land},
\Elim{\lif},
\Elim{\lnot} અને
\Elim{\lforall}નું
રૂપાંતર પરિણામની
\Log{G2i}
!!{proof}માં
\Cut{} અનુમાનો
આપે છે.
જોકે શરૂઆતની
\Log{N2i}-!!{proof}
સામાન્ય સ્વરૂપમાં
હોય, તો આ
ટાળી શકાય છે.
\sourcecorrection{OLMD-263}{મૂળનું છેલ્લું વાક્ય ``If the N2i-proof we start from, ...'' અધૂરું છે; આગળનું પ્રતિજ્ઞાપન સ્પષ્ટ કરે છે કે કટ ટાળવા શરૂઆતની સાબિતી સામાન્ય સ્વરૂપમાં હોવી જોઈએ.}

સિક્વન્ટોની
આવૃત્તિઓની
યાદી
$S_1$, \dots,~$S_n$
ને~\Log{N2i}-!!{proof}~$\delta$માં
\emph{શાખા} કહીએ,
જો~$S_1$
સ્વયંસિદ્ધ હોય
અને~$S_i$
કોઈ અનુમાનનું
મુખ્ય આધારવિધાન
હોય ત્યારે
$S_{i+1}$
તેનું નિષ્કર્ષ
હોય. બીજા
શબ્દોમાં, શાખા
સ્વયંસિદ્ધોથી
નીચે તરફ
$\delta$માં
સિક્વન્ટો અનુસરે
છે, પણ
!!{elimination}
નિયમના ગૌણ
આધારવિધાન પર
અટકે છે.
$S_n$ અંતિમ
સિક્વન્ટ હોય
એવી~$\delta$ની
શાખાને
\emph{મુખ્ય શાખા}
કહીએ. દરેક
!!{proof}~$\delta$માં
ઓછામાં ઓછી
એક મુખ્ય શાખા
હોય છે, કારણ
કે દરેક નિયમને
ઓછામાં ઓછું
એક મુખ્ય
આધારવિધાન છે
(ફક્ત~\Intro{\land}ને
બે છે).


\begin{prop}\ollabel{prop:normal-N2i-to-G2i}
જો~$\delta$ એ
$\Gamma \Sequent !A$ની
સામાન્ય
$\Log{N2c}$ અથવા
$\Log{N2i}$-!!{proof}
હોય, તો
$\Gamma' \Sequent \Delta$ની
(કટ-રહિત)
$\Log{G2c}$-!!{proof}
(અથવા
$\Log{G2i}$-!!{proof})~$\pi$
હોય છે. અહીં
$\Gamma$માંથી
ચિહ્નો દૂર કરતાં
મળતો !!{formula}sનો
બહુગણ~$\Gamma'$
છે; અને~$!A$
એ~$\lfalse$
ન હોય તો
$\Delta = \{!A\}$,
જ્યારે હોય તો
$\Delta = \emptyset$.
\end{prop}

\begin{proof}
આ વખતે
$\Gamma \Sequent !A$ની
!!{proof}~$\delta$માં
અનુમાનોની સંખ્યા
પર આગમન કરીએ.
$\delta$માં
એકેય અનુમાન
ન હોય, તો
તે સ્વયંસિદ્ધ
$x:!A \Sequent !A$
છે અને~$\pi$
અનુરૂપ સ્વયંસિદ્ધ
$!A \Sequent !A$
છે; અથવા
$!A
\ident \lfalse$
હોય તો
$\lfalse \Sequent \ $
છે.

જો~$\delta$નો
અંત અનુમાન-નિયમમાં
થાય, તો કિસ્સા
જુદા પાડીએ:
\begin{enumerate}
  \item $R$ એ
  !!a{introduction}
  નિયમ અથવા~\FalseInt
  હોય, તો
  રૂપાંતર
  \olref[nat][tni]{prop:N2-to-G2}ની
  સાબિતીની જેમ
  જ થાય છે
  અને~\Cut
  જરૂરી નથી.

  $\delta$નો
  છેલ્લો નિયમ~$R$
  !!a{elimination}
  નિયમ હોય,
  તો નીચેની
  વાતો નોંધો:
  \begin{enumerate}
    \item $\delta$ની
    મુખ્ય શાખા
    કોઈ !!{introduction}
    નિયમમાંથી
    પસાર થતી
    નથી. નહિતર
    એ શાખા પર
    !!a{introduction}
    નિયમ અથવા
    \FalseInt{}
    પછી
    !!a{elimination}
    નિયમ આવત
    અને~$\delta$
    સામાન્ય ન
    રહેત.
    \item $\delta$ની
    મુખ્ય શાખાઓ
    પર !!{introduction}
    નિયમો નથી,
    તેથી મુખ્ય
    શાખા એક
    જ છે.
    (માત્ર
    \Intro{\land}ને
    બે મુખ્ય
    આધારવિધાન છે;
    એટલે એકથી
    વધુ મુખ્ય
    શાખાઓ હોય
    તો તે
    \Intro{\land}ના
    આધારવિધાનોમાંથી
    પસાર થવી
    પડે.)
    \item મુખ્ય
    શાખામાં
    \Elim{\lor}
    અથવા
    \Elim{\lexists}નું
    મુખ્ય આધારવિધાન
    આવે, તો
    એવો ફક્ત
    એક જ
    નિયમ હોય
    અને તે
    છેલ્લો
    નિયમ~$R$
    હોય.
    (નહિતર
    મુખ્ય શાખામાં
    \Elim{\lor}
    કે
    \Elim{\lexists}
    એવો હોય
    કે તેનું
    નિષ્કર્ષ
    !!a{elimination}
    નિયમનું મુખ્ય
    આધારવિધાન
    બને; પછી
    ક્રમપરિવર્તન-રૂપાંતર
    લગાડી શકાય,
    એટલે સાબિતી
    સામાન્ય ન
    હોય.)
  \item \Elim{\lor}
  અને
  \Elim{\lexists}
  સિવાયના
  !!{elimination}
  નિયમો ધારણાઓ
  મુક્ત કરતા
  નથી, એટલે
  સિક્વન્ટના
  ડાબા સંદર્ભને
  બદલતા નથી.
  આ બંને નિયમો
  ફક્ત ગૌણ
  આધારવિધાનની
  ઉપરની ધારણાઓ
  મુક્ત કરે
  છે; એટલે
  ફક્ત પોતાના
  ગૌણ આધારવિધાનોનો
  ડાબો સંદર્ભ
  બદલે છે.
  બીજા શબ્દોમાં,
  $\delta$ની
  મુખ્ય શાખા
  પરના
  સિક્વન્ટ~$S_i$નો
  સંદર્ભ~$\Gamma_1$
  હોય, તો
  $S_i$ જે
  અનુમાનનું
  મુખ્ય આધારવિધાન
  છે તેના
  નિષ્કર્ષ
  $S_{i+1}$નો
  સંદર્ભ
  $\Gamma_1' \supseteq \Gamma_1$
  હોય છે.
  \item તેથી
  મુખ્ય શાખાનો
  સૌથી ઉપરનો
  સિક્વન્ટ~$S_1$
  જો
  $x: !C \Sequent !C$
  હોય, તો
  $x:!C \in \Gamma$.
  \end{enumerate}
  હવે~$!C$ના
  સ્વરૂપ પ્રમાણે
  કિસ્સા જુદા
  પાડીએ. મુખ્ય
  શાખાનું અંતિમ
  સિક્વન્ટ સિવાયનું
  દરેક~$S_i$
  !!a{elimination}
  નિયમનું મુખ્ય
  આધારવિધાન છે;
  તેથી આ વાત
  $x:!C \Sequent !C$
  માટે પણ
  સાચી છે.
  \sourcecorrection{OLMD-264}{મૂળ કહે છે કે મુખ્ય શાખાનો દરેક $S_i$ લોપ-નિયમનું મુખ્ય આધારવિધાન છે; અંતિમ સિક્વન્ટ પછી કોઈ અનુમાન જ નથી. દાવો અંતિમ સિવાયના $S_i$ માટે છે.}

  \item $!C \ident !D \land !E$
  હોય. ત્યારે~$\delta$નું
  સ્વરૂપ આ છે:
  \[
  \Axiom$x:!D \land !E \fCenter !D \land !E$
  \RightLabel{\Elim{\land}[1]}
  \UnaryInf$x:!D \land !E \fCenter !D$
  \Deduce$x:!D \land !E, \Gamma_1 \fCenter !A$
  \DisplayProof
  \]
  મુખ્ય શાખામાં
  $x:!D \land !E \Sequent !D$
  આવે છે. તેમાં
  \Intro{\lif}
  કે~\Intro{\lnot}નું
  મુખ્ય આધારવિધાન
  નથી, તેમજ
  \Elim{\lor}
  કે~\Elim{\lexists}નું
  ગૌણ આધારવિધાન
  પણ નથી; તેથી
  $x:!D
  \land !E$માંનાં
  સંદર્ભ-!!{formula}s
  આખી મુખ્ય
  શાખામાં યથાવત્
  રહે છે. તેથી
  અંતિમ સિક્વન્ટના
  સંદર્ભમાં
  $x:!D \land !E$
  હોવું જ પડે.
  વળી,
  !!{proof}~$\delta$માં
  \Elim{\land}માં
  પૂરી થતી
  ઉપ-!!{proof}ના
  સ્થાને
  $x:!D \Sequent !D$
  મૂકીને અને
  મુખ્ય શાખાના
  દરેક સિક્વન્ટના
  સંદર્ભમાં
  $x:!D \land !E$ના
  સ્થાને
  $x:!D$
  મૂકીને
  !!{proof}~$\delta_1$
  મેળવીએ; એટલે
  આને
  \[
  \bottomAlignProof
  \Axiom$x:!D \land !E \fCenter !D \land !E$
  \RightLabel{\Elim{\land}[1]}
  \UnaryInf$x:!D \land !E \fCenter !D$
  \Deduce$x:!D \land !E, \Gamma_1 \fCenter !A$
  \DisplayProof
  \quad\text{માં ફેરવો}\quad
  \bottomAlignProof
  \Axiom$x:!D \fCenter !D$
  \Deduce$x:!D, \Gamma_1 \fCenter !A$
  \DisplayProof
  \]
  આગમન-પરિકલ્પના
  $\delta_1$ને
  લાગુ પડે છે,
  કારણ કે~$\delta$માં
  અનુમાનોની સંખ્યા
  $\delta_1$માંની
  સંખ્યા કરતાં~$1$
  વધુ છે.
  \sourcecorrection{OLMD-265}{મૂળમાં આગમનથી બનેલી સાબિતીને $\delta_1$ કહીને પછી અનુમાનોની ગણતરીમાં અપરિભાષિત $\delta'$ વપરાયું છે; અહીં તે જ $\delta_1$ છે.}

  આગમન-પરિકલ્પનાથી
  $!D, \Gamma_1' \Sequent !A$ની
  \Log{G2i}-!!{proof}~$\pi_1$
  મળે છે.
  \LeftR{\land} લગાવી
  !!{proof}~$\pi$
  મેળવીએ:
  \sourcecorrection{OLMD-266}{મૂળ અહીં G2iની ઉપસાબિતીનું પૂર્વાંગ $x:!D, \Gamma_1$ લખે છે; G2iમાં $x:$ ચિહ્ન નથી અને સંદર્ભ $\Gamma_1'$ છે, જેમ નીચેના વૃક્ષમાં છે.}
  \[
  \AxiomC{}
  \RightLabel{$\pi_1$}
  \Deduce$!D, \Gamma_1' \fCenter !A$
  \RightLabel{\LeftR{\land}}
  \UnaryInf$!D \land !E, \Gamma_1' \fCenter !A$
  \DisplayProof
  \]
  $!A \ident \lfalse$
  હોય તો
  \RightR{\Weakening}
  ઉમેરીએ, જેમ કે:
  \[
  \AxiomC{}
  \RightLabel{$\pi_1$}
  \Deduce$!D, \Gamma_1' \fCenter $
  \RightLabel{\LeftR{\land}}
  \UnaryInf$!D \land !E, \Gamma_1' \fCenter $
  \RightLabel{\RightR{\Weakening}}
  \UnaryInf$!D \land !E, \Gamma_1' \fCenter !A$
  \DisplayProof
  \]
  \item જો
  $!C \ident !D \lor !E$
  હોય, તો તે
  \Elim{\lor}નું
  મુખ્ય આધારવિધાન
  હોવું જોઈએ
  અને આ અનુમાન
  મુખ્ય શાખાનું
  છેલ્લું અનુમાન~$R$
  હોવું જોઈએ.
  એટલે~$\delta$નું
  સ્વરૂપ આ છે:
  \[
  \Axiom$x:!D \lor !E \fCenter !D \lor !E$
  \AxiomC{}
  \RightLabel{$\delta_1$}
  \Deduce$y:!D, \Gamma_1 \fCenter!A$
  \AxiomC{}
  \RightLabel{$\delta_2$}
  \Deduce$z:!E, \Gamma_2 \fCenter !A$
  \RightLabel{\Elim{\lor}}
  \TrinaryInf$x:!D \lor !E, \Gamma_1, \Gamma_2 \fCenter !A$
  \DisplayProof
  \]
  આગમન-પરિકલ્પના
  !!{proof}s~$\delta_1$
  અને~$\delta_2$ને
  લાગુ પડે છે;
  તેથી
  $!D, \Gamma_1' \Sequent !A$
  અને
  $!E, \Gamma_2' \Sequent !A$ની
  \Log{G2i}-!!{proof}s
  $\pi_1$ અને
  $\pi_2$ મળે છે.
  $\LeftR{\lor}$
  લગાવી~$\pi$
  મેળવીએ:
  \[
  \AxiomC{}
  \RightLabel{$\pi_1$}
  \Deduce$!D, \Gamma_1' \fCenter !A$
  \AxiomC{}
  \RightLabel{$\pi_2$}
  \Deduce$!E, \Gamma_2' \fCenter !A$
  \RightLabel{\LeftR{\lor}}
  \BinaryInf$!D \lor !E, \Gamma_1', \Gamma_2' \fCenter !A$
  \DisplayProof
  \]
  $\Elim{\lor}$
  $y: !D$ અથવા
  $z: !E$ને
  મુક્ત ન કરતું
  હોય (એટલે
  તે ગૌણ
  આધારવિધાનોના
  સંદર્ભમાં
  ન હોય), અથવા
  $!A$ એ
  $\lfalse$
  હોય, તો
  કેટલાક
  \LeftR{\Weakening}
  અને
  \RightR{\Weakening}
  ઉમેરવા પડે.
  દાખલા તરીકે:
  \[
  \AxiomC{}
  \RightLabel{$\pi_1$}
  \Deduce$\Gamma_1' \fCenter $
  \RightLabel{\LeftR{\Weakening}}
  \UnaryInf$!D, \Gamma_1' \fCenter $
  \AxiomC{}
  \RightLabel{$\pi_2$}
  \Deduce$\Gamma_2' \fCenter $
  \RightLabel{\LeftR{\Weakening}}
  \UnaryInf$!E, \Gamma_2' \fCenter $
  \RightLabel{\LeftR{\lor}}
  \BinaryInf$!D \lor !E, \Gamma_1', \Gamma_2' \fCenter $
  \RightLabel{\RightR{\Weakening}}
  \UnaryInf$!D \lor !E, \Gamma_1', \Gamma_2' \fCenter !A$
  \DisplayProof
  \]
  
  \item $!C \ident !D \lif !E$
  હોય. ત્યારે~$\delta$નું
  સ્વરૂપ આ છે:
  \[
  \Axiom$x:!D \lif !E \fCenter !D \lif !E$
  \AxiomC{}
  \RightLabel{$\delta_1$}
  \Deduce$\Gamma_1 \fCenter!D$
  \RightLabel{\Elim{\lif}}
  \BinaryInf$x:!D \lif !E, \Gamma_1 \fCenter !E$
  \RightLabel{$\delta_2$}
  \Deduce$x:!D \lif !E, \Gamma_1, \Gamma_2 \fCenter !A$
  \DisplayProof
  \]
  અહીં નોંધો કે
  $x:!D \lif !E,
  \Gamma_1 \Sequent !E$
  ધરાવતી મુખ્ય
  શાખામાં
  \Intro{\lif}
  કે~\Intro{\lnot}નું
  મુખ્ય આધારવિધાન
  નથી, તેમજ
  \Elim{\lor}
  કે~\Elim{\lexists}નું
  ગૌણ આધારવિધાન
  પણ નથી.
  તેથી
  $x:!D \lif !E, \Gamma_1$ના
  સંદર્ભ-!!{formula}s
  મુખ્ય શાખામાં
  યથાવત્ રહે
  છે. આથી
  અંતિમ સિક્વન્ટના
  સંદર્ભમાં
  $x:!D \lif !E, \Gamma_1$
  આવવું જ જોઈએ.
  વળી,
  !!{proof}-ખંડ~$\delta_2$માં
  \Elim{\lif}માં
  પૂરી થતી
  ઉપ-!!{proof}ના
  સ્થાને
  $x:!E \Sequent !E$
  મૂકીએ અને
  મુખ્ય શાખાના
  દરેક સિક્વન્ટના
  સંદર્ભમાં
  $x:!D \lif !E,
  \Gamma_1$ના
  સ્થાને~$x:!E$
  મૂકીએ, તો
  યોગ્ય
  !!{proof}~$\delta_2'$
  મળે છે;
  એટલે આને
  \[
  \bottomAlignProof
  \Axiom$x:!D \lif !E, \Gamma_1 \fCenter !E$
  \RightLabel{$\delta_2$}
  \Deduce$x:!D \lif !E, \Gamma_1, \Gamma_2 \fCenter !A$
  \DisplayProof
  \quad\text{માં ફેરવો}\quad
  \bottomAlignProof
  \Axiom$x:!E \fCenter !E$
  \RightLabel{$\delta_2'$}
  \Deduce$x:!E, \Gamma_2 \fCenter !A$
  \DisplayProof
  \]
  આગમન-પરિકલ્પના
  $\delta_1$ અને
  $\delta_2'$ને
  લાગુ પડે છે,
  કારણ કે~$\delta$માં
  અનુમાનોની સંખ્યા
  $\delta_1$ અને
  $\delta_2$માંની
  અનુમાનોની
  સંખ્યાના સરવાળાથી~$1$
  વધુ છે; આ
  એક અનુમાન
  મુખ્ય શાખાની
  શરૂઆતનું
  \Elim{\lif}
  છે. (જો
  આ જ
  $\delta$નું
  છેલ્લું
  અનુમાન~$R$
  હોય, તો
  $\delta_1$માં
  $0$ અનુમાનો
  છે.)

  આગમન-પરિકલ્પનાથી
  $\Gamma_1' \Sequent !D$ની
  \Log{G2i}-!!{proof}~$\pi_1$
  અને
  $!E, \Gamma_2' \Sequent !A$ની
  સાબિતી~$\pi_2$
  મળે છે.
  \LeftR{\lif} લગાવી
  !!{proof}~$\pi$
  મેળવીએ:
  \sourcecorrection{OLMD-267}{મૂળના આગમન-પરિણામમાં $\pi_1$ માટે $\Gamma_1$નું ચિહ્નવિહોણું સ્વરૂપ દર્શાવતું પ્રાઇમ નથી અને $\pi_2$ માટે $\Gamma_2'$ને બદલે $\Gamma_1'$ લખાયું છે; નીચેનું ડાબા અનુસરણનું વૃક્ષ યોગ્ય સંદર્ભો બતાવે છે.}
  \[
  \AxiomC{}
  \RightLabel{$\pi_1$}
  \Deduce$\Gamma_1' \fCenter !D$
  \AxiomC{}
  \RightLabel{$\pi_2$}
  \Deduce$!E, \Gamma_2' \fCenter !A$
  \RightLabel{\LeftR{\lif}}
  \BinaryInf$!D \lif !E, \Gamma_1', \Gamma_2' \fCenter !A$
  \DisplayProof
  \]
  $!D \ident \lfalse$
  અને/અથવા
  $!A \ident \lfalse$
  હોય, તો
  એક કે બે
  \RightR{\Weakening}
  ઉમેરીએ,
  જેમ કે:
  \sourcecorrection{OLMD-268}{મૂળમાં શરતનો પહેલો ભાગ માત્ર $!D$ લખાયો છે; ખાલી ઉત્તરાંગનો કિસ્સો $!D \ident \lfalse$ છે, જેમ તરત નીચેની જમણી શિથિલીકરણની કડી બતાવે છે.}
  \[
  \AxiomC{}
  \RightLabel{$\pi_1$}
  \Deduce$\Gamma_1' \fCenter $
  \RightLabel{\RightR{\Weakening}}
  \UnaryInf$\Gamma_1' \fCenter !D$
  \AxiomC{}
  \RightLabel{$\pi_2$}
  \Deduce$!E, \Gamma_2' \fCenter $
  \RightLabel{\RightR{\Weakening}}
  \UnaryInf$!E, \Gamma_2' \fCenter !A$
  \RightLabel{\LeftR{\lif}}
  \BinaryInf$!D \lif !E, \Gamma_1', \Gamma_2' \fCenter !A$
  \DisplayProof
  \]

\end{enumerate}
બાકીના કિસ્સા
આવી જ રીતે
સંભાળી શકાય
છે અને તેમને
કસરત તરીકે
મૂક્યા છે.
\end{proof}

\begin{cor}
સામાન્ય
\Log{N1i} અથવા
\Log{N2i}-!!{proof}માં
આવતું દરેક
!!{formula}
અંતિમ
!!{formula} અથવા
અંતિમ સિક્વન્ટનું
ઉપ-!!{formula}
છે.
\end{cor}

\begin{proof}
  \olref{prop:normal-N2i-to-G2i}
  પ્રમાણે દરેક
  સામાન્ય
  \Log{N2i}-!!{proof}ને
  \Log{G2i}-!!{proof}માં
  રૂપાંતરિત કરી
  શકાય છે.
  સાબિતી જોવાથી
  દેખાય છે કે
  \Log{N2i}-!!{proof}માં
  આવતું દરેક
  !!{formula}
  \Log{G2i}-!!{proof}માં
  પણ આવે છે.
  \Log{G2i}માં
  \Cut{} નિયમ
  નથી, તેથી
  તેમાં
  ઉપ-!!{formula}
  ગુણધર્મ છે.
\end{proof}

\end{document}
